\chapter{Permissions, seen rather than described}
\label{ch:roles}

Permission systems are usually taught as a grid, and the grid is forgettable.
This chapter teaches it differently: the same page, photographed three times,
signed in as three different people.

\section{One page, three people}

\shot{11-incidents-as-viewer.png}{\textbf{Sofia Rossi --- viewer.} She can read
everything and change nothing. No ``Report incident''. No ``Export CSV''.}

\shot{05-incidents-list.png}{\textbf{Daniel Okoye --- incident commander.} The
same data, plus the ability to act on it.}

\shot{12-incidents-as-admin.png}{\textbf{Riya Sharma --- administrator.} As the
commander, plus organisation-level control and an ``Admin'' item in the
navigation.}

Nothing about the data changed between those three pictures. What changed is
what is \emph{offered}.

\section{Where the difference is defined}

Exactly one file, and it is short enough to read in full.

\where{api/app/Enums/OrganizationRole.php}
\begin{lstlisting}[language=PHP]
public function permissions(): array
{
    $viewer = [
        Permission::IncidentView,
        Permission::PostmortemView,
        Permission::ServiceView,
        Permission::MemberView,
        Permission::MetricsView,
    ];

    $reporter = [
        ...$viewer,
        Permission::IncidentCreate,
        Permission::CommentCreate,
    ];

    $responder = [
        ...$reporter,
        Permission::IncidentUpdate,
        Permission::IncidentTransition,
        Permission::UpdateCreate,
    ];

    $commander = [
        ...$responder,
        Permission::IncidentAssign,
        Permission::IncidentCommand,
        Permission::CommentModerate,
        Permission::PostmortemManage,
        Permission::PostmortemPublish,
        Permission::ExportRun,
    ];

    return match ($this) {
        self::Viewer            => $viewer,
        self::Reporter          => $reporter,
        self::Responder         => $responder,
        self::IncidentCommander => $commander,
        self::Administrator     => [
            ...$commander,
            Permission::IncidentDelete,
            Permission::ServiceManage,
            Permission::MemberManage,
            Permission::AuditView,
            Permission::OrganizationManage,
        ],
    };
}
\end{lstlisting}

\subsection{Read the shape, not just the contents}

The \fpath{...} is PHP's spread operator: \fpath{...\$viewer} means ``everything
a viewer can do, and then these as well''. Each role is literally written as the
previous one plus a few additions.

\note{This makes the hierarchy \emph{structural} rather than a convention
somebody has to maintain. It is impossible to give a reporter a permission a
viewer lacks by accident, because a reporter's list is built from the viewer's.
A flat table --- five arrays written out in full --- would allow that mistake,
and would allow the five lists to drift apart over years of small edits.}

\subsection{Permissions are an enum, not strings}

\fpath{Permission::IncidentCreate}, not \fpath{'incident.create'}. Typing
\fpath{Permission::IncidentCreat} is an error your editor catches before you
save; typing \fpath{'incident.creat'} is a permission check that silently returns
false forever --- and a silently-false permission check fails \emph{open} or
\emph{closed} depending on which side of the comparison it lands, neither of
which you will notice in testing.

\section{Now the same picture from the other side}

The screenshots showed the frontend. Here is what actually enforces it:

\where{api/app/Enums/OrganizationRole.php}
\begin{lstlisting}[language=PHP]
public function has(Permission $permission): bool
{
    return in_array($permission, $this->permissions(), strict: true);
}
\end{lstlisting}

\fpath{strict: true} compares type as well as value. Without it PHP's loose
comparison can produce surprising equalities, and a permission check is the last
place you want a surprise.

\warnbox{Say this to yourself until it is automatic: \textbf{the missing button
is not the security}. A viewer who opens the browser console and sends the
request by hand meets this check on the server, and is refused with 403. If you
ever find yourself relying on a hidden control to prevent something, you have
not prevented it.}

\section{The full grid, for reference}

\begin{center}
\begin{tabular}{@{}>{\raggedright\arraybackslash}p{5.2cm}ccccc@{}}
\toprule
\textbf{Can\ldots} & \rotatebox{60}{viewer} & \rotatebox{60}{reporter} & \rotatebox{60}{responder} & \rotatebox{60}{commander} & \rotatebox{60}{admin} \\
\midrule
view incidents, services, metrics & yes & yes & yes & yes & yes \\
create an incident & --- & yes & yes & yes & yes \\
comment & --- & yes & yes & yes & yes \\
change status & --- & --- & yes & yes & yes \\
post an update & --- & --- & yes & yes & yes \\
assign people & --- & --- & --- & yes & yes \\
change severity & --- & --- & --- & yes & yes \\
publish a postmortem & --- & --- & --- & yes & yes \\
export CSV & --- & --- & --- & yes & yes \\
manage services & --- & --- & --- & --- & yes \\
manage members & --- & --- & --- & --- & yes \\
read the audit log & --- & --- & --- & --- & yes \\
\bottomrule
\end{tabular}
\end{center}

\section{Try it yourself}

With the stack running, sign in as each of these in turn. The password for every
demo account is \fpath{incidentflow}.

\begin{itemize}[leftmargin=1.4cm]
\cbox \fpath{viewer@incidentflow.test} --- confirm ``Report incident'' is absent.
\cbox \fpath{reporter@incidentflow.test} --- it appears. Open an incident and
  confirm you cannot change its status.
\cbox \fpath{responder@incidentflow.test} --- now you can.
\cbox \fpath{commander@incidentflow.test} --- assignment and severity appear.
\cbox \fpath{admin@incidentflow.test} --- an ``Admin'' item appears in the
  navigation.
\end{itemize}

\expect{Five sign-ins, five visibly different versions of the same application.
If two of them look identical, you are still signed in as the previous user ---
sign out properly between them.}
